\appendix
\section{Ansible}
\label{app:Ansible}
\subsubsection{ansible.cfg}
\label{subsec:ansible}
\begin{minted}{YAML}
Balls
\end{minted}


\section{Ansible Roles}
\label{app:Ansible-roles}

\subsection{iptables.yml}
\subsubsection{tasks/main.yml}
\begin{minted}{YAML}
Balls
\end{minted}

\subsubsection{templates/class-a.j2}
\begin{minted}{Python}
Class A
\end{minted}

\subsubsection{templates/class-b.j2}
\begin{minted}{Python}
Class B
\end{minted}

\subsubsection{templates/class-c.j2}
\begin{minted}{Python}
Class C
\end{minted}


\section{Firewall Configuration Files}
\label{app:config}
% Include complete configuration files

\section{Firewall Rulesets}
\label{app:rules}

\subsection{Class-A Network Rules (WAN)}
\label{subsec:class-a-rules}

The Class-A network rules apply to the WAN interface (eth0) connecting to the external internet. These rules implement strict ingress filtering, allow essential services, enable packet forwarding to internal networks, and include anti-spoofing measures.

\begin{verbatim}
*filter
:INPUT DROP [0:0]
:FORWARD DROP [0:0]
:OUTPUT ACCEPT [0:0]

-A INPUT -i lo -j ACCEPT
-A INPUT -m state --state ESTABLISHED,RELATED -j ACCEPT

-A INPUT -p tcp --dport 22 -m state --state NEW -j ACCEPT
-A INPUT -i eth0 -p icmp --icmp-type echo-request -j ACCEPT

-A INPUT -i eth0 -p tcp --dport 80 -j ACCEPT
-A INPUT -i eth0 -p tcp --dport 443 -j ACCEPT

-A INPUT -i eth0 -p udp --dport 53 -j ACCEPT
-A INPUT -i eth0 -p tcp --dport 53 -j ACCEPT

-A INPUT -j LOG --log-prefix "INPUT-DROP: " --log-level 4
-A INPUT -j DROP

-A FORWARD -m state --state ESTABLISHED,RELATED -j ACCEPT

-A FORWARD -i eth1 -o eth0 -j ACCEPT
-A FORWARD -i eth2 -o eth0 -j ACCEPT

-A FORWARD -i eth1 -o eth2 -j DROP
-A FORWARD -i eth2 -o eth1 -j DROP

-A FORWARD -p tcp --dport 22 -j DROP
-A FORWARD -p tcp --dport 23 -j DROP
-A FORWARD -p tcp --dport 3389 -j DROP

-A FORWARD -p tcp --tcp-flags ALL NONE -j DROP
-A FORWARD -p tcp --tcp-flags ALL ALL -j DROP

-A FORWARD -j LOG --log-prefix "FORWARD-DROP: " --log-level 4
-A FORWARD -j DROP

COMMIT

*nat
:PREROUTING ACCEPT [0:0]
:INPUT ACCEPT [0:0]
:OUTPUT ACCEPT [0:0]
:POSTROUTING ACCEPT [0:0]

-A POSTROUTING -o eth0 -j MASQUERADE

-A POSTROUTING -s 172.16.0.0/16 -o eth0 -j MASQUERADE
-A POSTROUTING -s 192.168.0.0/24 -o eth0 -j MASQUERADE

COMMIT

*mangle
:PREROUTING ACCEPT [0:0]
:INPUT ACCEPT [0:0]
:FORWARD ACCEPT [0:0]
:OUTPUT ACCEPT [0:0]
:POSTROUTING ACCEPT [0:0]

-A PREROUTING -m conntrack --ctstate INVALID -j DROP

-A PREROUTING -p tcp ! --syn -m conntrack --ctstate NEW -j DROP

-A PREROUTING -p tcp --tcp-flags ALL NONE -j DROP
-A PREROUTING -p tcp --tcp-flags ALL ALL -j DROP
-A PREROUTING -p tcp --tcp-flags SYN,FIN SYN,FIN -j DROP
-A PREROUTING -p tcp --tcp-flags SYN,RST SYN,RST -j DROP

-A PREROUTING -f -j DROP

COMMIT
\end{verbatim}

\subsection{Class-B Network Rules (LAN)}
\label{subsec:class-b-rules}

The Class-B network rules apply to the LAN interface (eth1) serving the trusted network segment. These rules permit internal management traffic, DNS/DHCP services, and restrict forbidden protocols while maintaining isolation from LAN2.

\begin{verbatim}
*filter
:INPUT DROP [0:0]
:FORWARD DROP [0:0]
:OUTPUT ACCEPT [0:0]

-A INPUT -i lo -j ACCEPT
-A INPUT -m state --state ESTABLISHED,RELATED -j ACCEPT

-A INPUT -p tcp --dport 22 -m state --state NEW -j ACCEPT
-A INPUT -i eth1 -p icmp --icmp-type echo-request -j ACCEPT

-A INPUT -i eth1 -p udp --dport 67 -j ACCEPT
-A INPUT -i eth1 -p udp --dport 68 -j ACCEPT

-A INPUT -i eth1 -p udp --dport 53 -j ACCEPT
-A INPUT -i eth1 -p tcp --dport 53 -j ACCEPT

-A INPUT -i eth1 -p tcp --dport 80 -j ACCEPT
-A INPUT -i eth1 -p tcp --dport 443 -j ACCEPT

-A INPUT -i eth1 -p tcp --dport 3000 -j ACCEPT
-A INPUT -i eth1 -p tcp --dport 8080 -j ACCEPT

-A INPUT -s 172.16.0.0/16 -i eth1 -j ACCEPT

-A INPUT -j LOG --log-prefix "INPUT-DROP-LAN: " --log-level 4
-A INPUT -j DROP

-A FORWARD -m state --state ESTABLISHED,RELATED -j ACCEPT

-A FORWARD -i eth1 -o eth0 -s 172.16.0.0/16 -j ACCEPT

-A FORWARD -i eth1 -o eth2 -j DROP

-A FORWARD -i eth0 -o eth1 -d 172.16.0.0/16 -m state --state ESTABLISHED,RELATED -j ACCEPT

-A FORWARD -p tcp --dport 25 -j DROP
-A FORWARD -p tcp --dport 135 -j DROP
-A FORWARD -p tcp --dport 139 -j DROP
-A FORWARD -p tcp --dport 445 -j DROP

-A FORWARD -j LOG --log-prefix "FORWARD-DROP-LAN: " --log-level 4
-A FORWARD -j DROP

COMMIT

*nat
:PREROUTING ACCEPT [0:0]
:INPUT ACCEPT [0:0]
:OUTPUT ACCEPT [0:0]
:POSTROUTING ACCEPT [0:0]

-A POSTROUTING -s 172.16.0.0/16 -o eth0 -j MASQUERADE

COMMIT

*mangle
:PREROUTING ACCEPT [0:0]
:INPUT ACCEPT [0:0]
:FORWARD ACCEPT [0:0]
:OUTPUT ACCEPT [0:0]
:POSTROUTING ACCEPT [0:0]

-A PREROUTING -i eth1 -m conntrack --ctstate INVALID -j DROP

-A PREROUTING -i eth1 -p tcp ! --syn -m conntrack --ctstate NEW -j DROP

-A PREROUTING -i eth1 -s 172.16.0.0/16 -j ACCEPT

COMMIT
\end{verbatim}

\subsection{Class-C Network Rules (LAN2)}
\label{subsec:class-c-rules}

The Class-C network rules apply to the LAN2 interface (eth2) serving untrusted devices. These rules implement the strictest filtering with rate limiting, block access to internal networks, restrict dangerous services (SSH, RDP, database ports), and implement anti-spoofing measures to prevent devices claiming RFC1918 addresses.

\begin{verbatim}
*filter
:INPUT DROP [0:0]
:FORWARD DROP [0:0]
:OUTPUT ACCEPT [0:0]

-A INPUT -i lo -j ACCEPT
-A INPUT -m state --state ESTABLISHED,RELATED -j ACCEPT

-A INPUT -i eth2 -p icmp --icmp-type echo-request -j ACCEPT

-A INPUT -i eth2 -p udp --dport 67 -j ACCEPT
-A INPUT -i eth2 -p udp --dport 68 -j ACCEPT

-A INPUT -i eth2 -p udp --dport 53 -j ACCEPT
-A INPUT -i eth2 -p tcp --dport 53 -j ACCEPT

-A INPUT -i eth2 -p tcp --dport 80 -j ACCEPT
-A INPUT -i eth2 -p tcp --dport 443 -j ACCEPT

-A INPUT -i eth2 -p tcp --dport 22 -m state --state NEW -j ACCEPT

-A INPUT -i eth2 -s 192.168.0.0/24 -m limit --limit 10/min -j ACCEPT

-A INPUT -j LOG --log-prefix "INPUT-DROP-GUEST: " --log-level 4
-A INPUT -j DROP

-A FORWARD -m state --state ESTABLISHED,RELATED -j ACCEPT

-A FORWARD -i eth2 -o eth0 -s 192.168.0.0/24 -j ACCEPT

-A FORWARD -i eth2 -o eth1 -j DROP

-A FORWARD -i eth2 -d 172.16.0.0/16 -j DROP

-A FORWARD -i eth0 -o eth2 -d 192.168.0.0/24 -m state --state ESTABLISHED,RELATED -j ACCEPT

-A FORWARD -i eth2 -p tcp --dport 22 -j DROP
-A FORWARD -i eth2 -p tcp --dport 23 -j DROP
-A FORWARD -i eth2 -p tcp --dport 3389 -j DROP

-A FORWARD -i eth2 -p tcp --dport 135 -j DROP
-A FORWARD -i eth2 -p tcp --dport 139 -j DROP
-A FORWARD -i eth2 -p tcp --dport 445 -j DROP

-A FORWARD -i eth2 -p tcp --dport 1433 -j DROP
-A FORWARD -i eth2 -p tcp --dport 3306 -j DROP
-A FORWARD -i eth2 -p tcp --dport 5432 -j DROP

-A FORWARD -i eth2 -m limit --limit 100/sec --limit-burst 200 -j ACCEPT

-A FORWARD -j LOG --log-prefix "FORWARD-DROP-GUEST: " --log-level 4
-A FORWARD -j DROP

COMMIT

*nat
:PREROUTING ACCEPT [0:0]
:INPUT ACCEPT [0:0]
:OUTPUT ACCEPT [0:0]
:POSTROUTING ACCEPT [0:0]

-A POSTROUTING -s 192.168.0.0/24 -o eth0 -j MASQUERADE

COMMIT

*mangle
:PREROUTING ACCEPT [0:0]
:INPUT ACCEPT [0:0]
:FORWARD ACCEPT [0:0]
:OUTPUT ACCEPT [0:0]
:POSTROUTING ACCEPT [0:0]

-A PREROUTING -i eth2 -m conntrack --ctstate INVALID -j DROP

-A PREROUTING -i eth2 -p tcp ! --syn -m conntrack --ctstate NEW -j DROP

-A PREROUTING -i eth2 -p tcp --tcp-flags ALL NONE -j DROP
-A PREROUTING -i eth2 -p tcp --tcp-flags ALL ALL -j DROP
-A PREROUTING -i eth2 -p tcp --tcp-flags SYN,FIN SYN,FIN -j DROP
-A PREROUTING -i eth2 -p tcp --tcp-flags SYN,RST SYN,RST -j DROP

-A PREROUTING -i eth2 -s 192.168.0.0/24 -j ACCEPT

-A PREROUTING -i eth2 -s 10.0.0.0/8 -j DROP
-A PREROUTING -i eth2 -s 172.16.0.0/12 -j DROP

COMMIT
\end{verbatim}

\section{Testing Scripts}
\label{app:scripts}
% Network testing and monitoring scripts

\section{Additional Screenshots}
\label{app:screenshots}
% Supporting visual evidence

\section{Command Outputs}
\label{app:outputs}
% Full command execution results